\uuid{H6oe}
\exo7id{1783}
\titre{exo7 1783}
\auteur{gineste}
\organisation{exo7}
\datecreate{2001-11-01}
\isIndication{false}
\isCorrection{false}
\chapitre{Topologie}
\sousChapitre{Autre}
\cpge{mpsi-mp2i.nombres-reels-suites.f, pcsi.nombres-reels-suites.e, ptsi.nombres-reels-suites.e, mpsi-mp2i.nombres-reels-suites.d, pcsi.nombres-reels-suites.c, ptsi.nombres-reels-suites.c, mp-mpi.topologie-espaces-normes.b, pc.espaces-vectoriels-normes.b, psi.espaces-vectoriels-normes.b, mp-mpi.topologie-espaces-normes.d, pc.espaces-vectoriels-normes.d, psi.espaces-vectoriels-normes.d, pt.fonctions-deux-trois-variables.a, mp-mpi.topologie-espaces-normes.g}
\module{Analyse}
\niveau{L2}
\difficulte{}

\contenu{
\question{
Soit $(u_n)_{n \geq 1}$ une suite r\'eelle. $\forall n \geq 1$, on pose $A_n=\{ u_p / \; p\geq n \}.$ D\'emontrer que l'ensemble des valeurs d'adh\'erence de la suite  $(u_n)_{n \geq 1}$ est $V= \bigcap_{n \geq 1} \overline{A_n}$, et qu'ainsi $V$ est ferm\'e. En d\'eduire que si la suite est born\'ee, alors l'ensemble $V$ est un compact non vide.
}
}
